\worldchapter{Horror}{horror}
\label{ch:horror}

\begin{multicols}{2}

The Phase Six Horror Expansion adds classic horror elements to the game. Characters can load silver ammunition into their weapons and encounter creatures from the darkest imaginations, as well as obscure items.

In addition, characters now have a potential stress level. If they encounter the non-worldly too intensely, there is a risk that they will lose control or even acquire a permanent quirk.

\section{Stress}

The character has a stress level that starts at 0 and a maximum level that they can withstand. This value is set to 10 by default, but can be adjusted using templates or similar methods.

Characters can gain stress when they encounter otherworldly entities or experience other abnormal events. Each creature is listed alongside the amount of stress that encountering it causes. This information consists of two values separated by a slash.

\begin{quote}
Example: A Spectre causes 1/2 stress when encountered.
\end{quote}

When a character encounters a creature that causes stress, they must perform a stress test. If they succeed, they receive as much stress as indicated before the slash. If the roll fails, the value after the slash applies.

If a character encounters multiple creatures at the same time, only one stress test is performed and stress is recorded once per character. However, special situations, such as 100 zombies approaching a character, may require separate stress tests, as determined by the game master.

\subsection{Base Stress}

A newly created character starts with a base stress level of 0, which indicates the amount of stress accumulated during adventures and the extent to which the character's mind has been affected. Base stress has no significance in the game; it simply represents the minimum value below which stress cannot fall.

Base stress can only be reduced in exceptional cases and through therapy. It can increase if a dread is resolved during a rest (see Dread).

If the base stress exceeds the maximum stress level, the character is permanently \textit{overcome by dread}. This means that the \textit{Resolve Dread} rule must be taken into account at each rest, and the character will continue to be affected by \textit{Dread} afterwards.

\subsection{Stress Test}

A stress test is performed whenever it is necessary to assess whether a character can withstand a stressful situation. This involves adding together the character's \textit{Logic} and \textit{Willpower} values, after which the corresponding number of dice are rolled. If the roll shows at least one success, the test is passed. The minimum roll corresponds to the character's minimum roll and is usually not altered.

\subsection{Reducing stress}

In order to reduce stress, the character must calm down and avoid encounters with otherworldly beings. This could involve taking a quiet moment or performing an activity within the game. Resting also reduces stress.

There are various ways to reduce stress in the game. What they all have in common is that they take more than an hour.

Some examples are:
\begin{itemize}
    \item Finding peace/meditation: -1 stress
    \item Writing in a diary: -1 stress
    \item Talking to someone: -2 stress
    \item petting a cat: -2 stress
    \item Taking certain medications/drugs: -X stress
\end{itemize}

During rest (see \cref{ch:wounds}), stress is reduced to the minimum possible (base stress or 0).

\section{Dread}

When a character's stress level is at its maximum, they are overcome by dread. This is a fixed state, with no different levels. Once a character is overcome by dread, their stress level remains at its maximum.

There is always a spontaneous effect when a character is overcome by dread. This effect lasts for a minimum of one hour and up to D6-Resistance hours. To determine the effect, 3d6 is rolled. The result is listed below:
\begin{itemize}
    \item 3–4: nausea and dizziness
    \item 5–8: anxiety
    \item 9–10: shock
    \item 11–12: panic
    \item 13–14: confusion
    \item 15–16: hallucinations
    \item 17–18: blindness
\end{itemize}

In addition to causing stress, creatures may have an ability that directly causes dread.

\subsection{Resolving dread}

Dread usually resolves during rest.

To achieve this, a stress test is performed during rest. If the test is successful, the \textit{Consumed by dread} condition is removed. If the test fails, the base stress increases by 1 and the character receives a quirk. In any case, the condition is removed.

\subsection{Dread and further stress}

If a character experiences further stress while in a state of dread, their stress level remains unchanged as it is already at its maximum value. Instead, the dice table for the effects of dread is rolled again for each new instance of stress, regardless of how high it is.

The character does not gain a quirk from this. A character can only be in the \textit{Consumed by dread} state once. This state can only be resolved by resting.

\section{Quirks}

A quirk is a permanent trait that a character develops in response to stress and fear. Every quirk has positive and negative aspects. Quirks can only be cured through lengthy therapy outside of the game.

\subsection{Acquiring quirks}

If a character is overcome by dread, the effect is resolved during a rest. This involves rolling a stress test. If the test fails, the character receives one base stress point and a quirk of their choice from the list of quirks.

This quirk is recorded on the character sheet, and its effects take effect immediately in the game. Quirks can have their own rules and can also alter the character's game values.

\subsection{Healing Quirks}

Quirks can only be healed outside of the game. This requires lengthy therapy. This must be done in consultation with the game master, for example when the story takes a break or when the player is playing a different character.

The corresponding quirk is then simply removed.
\end{multicols}
